\subsection{Cummulative Distribution Function (CDF/KVF)}
\begin{sectionbox}
    \begin{emphbox}
        Wahrscheinlichkeit, dass ein Ergebnis kleiner oder gleich $x$ ist
        \begin{equation*}
            F_X:\mathbb R \rightarrow [0,1], \; x\mapsto F_X(x)=P(\{X \le x\})
        \end{equation*}
        \begin{itemize}
            \item $X$: Experiment, von dem wir reden
            \item $x$: mögliches Ergebnis des Zufallsexperiments, hier die obere Schranke
            \item $F_X$ beschreibt das Experiment $X$ vollständig
        \end{itemize}
    \end{emphbox}

    \subsubsection{Eigenschaften}
    \begin{itemize}
        \item $F_X(x)$ monoton wachsend und rechtsseitig stetig: \\
        $\lim_{h\rightarrow 0^+}F_X(x + h)=F_X(x)$
        \item $F_X(x) \ge 0$
    \end{itemize}
\end{sectionbox}


\subsection{Diskrete Wahrscheinlichkeitsverteilungen}
$X:\Omega\rightarrow \mathbb N$ (oder $\mathbb Z, ...$)

\begin{sectionbox}
    \subsubsection{Probability Mass Function (PMF/WMF)}
    Wahrscheinlichkeit eines Ergebnisses $z$
    \begin{equation*}
        P_X(z)=P(\{X = z\})
    \end{equation*}
\end{sectionbox}

\begin{sectionbox}
    \subsubsection{CDF im diskreten Fall}
    \begin{equation*}
        F_X(x)=P(\{X \le x\}) = P(\{\{x_1,x_2,...,x_k\} | x_{1,...,k} \le x\})
    \end{equation*}
    \begin{equation*}
        =\sum_{x_i\le x}P(\{X=x_i\})=\sum_{x_i\le x} P_X(x_i)
    \end{equation*}
\end{sectionbox}

\subsection{Stetige Wahrscheinlichkeitsverteilungen}
\begin{emphbox}
    Wahrscheinlichkeiten von Intervallen
    \begin{equation*}
        P(\{a < X \le b\})=F_X(b) - F_X(a)
    \end{equation*}
\end{emphbox}

\subsubsection{Probability Density Function $f_X$ (PDF/WDF)}
\begin{sectionbox}
    \begin{emphbox}
        \begin{equation*}
            f_X(x)=\left.\dfrac{d}{dx}F_X(x)\right|_x \; \Longleftrightarrow \; F_X(x) = \int_{-\infty}^x f_X(x)\diff x
        \end{equation*}
    \end{emphbox}

    \textbf{Hauptintegralsatz}
    \begin{equation*}
        \int_a^bf_X(x)\diff x = F_X(b) - F_X(a) = P(\{a < X \le b\})
    \end{equation*}

\end{sectionbox}

\subsection{Mischverteilungen}
Problem z.B.: Diskretes Ereignis für $t = 0$ und stetige Ereignisse für $t > 0$. \\
\begin{sectionbox}
    z.B.: Wahrscheinlichkeit, dass Gerät bei Zeit $T=t$ kaputt geht.
    \begin{description}
        \item[Fall 1] $P(\{T=0\}) = p_0$ für defekt beim Einschalten, diskret (es wurde vor dem Einschalten schon kaputt)
        \item[Fall 2] Exponentielle Dichteverteilung für die Lebenszeit, vorausgesetzt das Einschalten wurde überlebt (bedingte Wahrscheinlichkeit), stetig: $f_T(t) = \lambda e^{-\lambda t} \; t \ge 0$
    \end{description}

    \begin{equation*}
        F_T(t) = \begin{cases}
            0 & t < 0 \\
            p_0 & t = 0 \\
            p_0 + (1-p_0)P(\{0 < T \le t\}) & t > 0
        \end{cases}
    \end{equation*}
    mit \begin{equation*}
        P(\{0 < T \le t\}) = \int_0^t f_T(\tau)\diff\tau
    \end{equation*}
\end{sectionbox}


\subsection{Mehrdimensionale Zufallsvariablen}
\begin{emphbox}
    \begin{equation*}
        F_{X_1X_2}(x_1, x_2) = P(\{X_1 \le x_1\}, \{X_2 \le x_2\})
    \end{equation*}
\end{emphbox}
\begin{sectionbox}
    \subsubsection{Stetige PDF}
    \begin{emphbox}
        \begin{equation*}
            f_{X_1X_2}(x_1,x_2)=\partial_{x_1}\partial_{x_2}\left(F_{X_1X_2}(x_1, x_2)\right)
        \end{equation*}
        $\Updownarrow$
        \begin{equation*}
            F_{X_1X_2}(x_1,x_2) = \int_{-\infty}^{x_1}\int_{-\infty}^{x_2}f_{X_1X_2}(\tau_1, \tau_2)\diff\tau_1\diff\tau_2
        \end{equation*}
    \end{emphbox}
\end{sectionbox}

\begin{sectionbox}
    \subsubsection{Marginalisierung}
    Einschränkung auf eine Zufallsvariable. Idee: \\
    \begin{equation*}
        \lim_{x_2\rightarrow\infty}F_{X_1X_2}(x_1, x_2) = P(\{X_1 \le x_1\})=F_{X_1}(x_1), \; x_2 \text{ bel.}
    \end{equation*}

    \begin{emphbox}
        \begin{equation*}
            f_{X_1}(x_1) = \int_{-\infty}^\infty f_{X_1X_2}(x_1,x_2)\diff x_2
        \end{equation*}
    \end{emphbox}
\end{sectionbox}